\documentclass{article}
\usepackage{/globals/de}
\usepackage{/globals/anki}
\ankiprefix{begriff-politik}
\ankitopic{Begriff Politik}
\ankishow
\begin{document}
\section{Begriff: Politik} 
Die Politik behandelt das Zusammenleben der Menschen. Sie wird von bestimmten \textquote{Gegenständen} ausgemacht; an erster Stelle unterschiedlichen gesellschaftlichen Problemen.

Oftmals geht es aber auch viel mehr um Macht.

Politik kann auch als Management von Interdependenzen gesehen werden.

Für die gesellschaftliche Verbindlichkeit ist Macht benötigt, an erster Stelle durch die Exekutive.
\ankinote{def}{Definition}{Zusammenleben Menschen, gesellschaftliche Probleme, Management von Interdependenzen}

\subsection{Das Politische System}
Das politische System ist der Ort an dem die Politik tatsächlisch stattfindet. Es beinhaltet die Institutionen und Regelungen in denen und mithilfe dessen Entscheidungen getroffen werden und Konflikte gelöst werden.

Bürger sind über Gruppen, Parteien \usw, Mitglieder des Systems; so entsteht Politik sehr dezentral.
\ankinote{system}{Politische System}{Ort der Politik, Institutionen, Bürger als Mitglieder, dezentral}

\subsection{Medien}
Die Medien werden oft auch die vierte Gewalt genannt. Sie haben die Macht Meinungen zu bilden, können Skandale aufdecken und wirken somit wie eine weitere Kontrollinstanz.
\ankinote{median}{Median}{Vierte Gewalt, Meinung, Skandale, Kontrollinstant}

\end{document}